\documentclass{article}
\usepackage[utf8]{inputenc}
\usepackage{helvet}
\usepackage{geometry}
\usepackage{xcolor}
\usepackage{eso-pic}
\usepackage{fancyhdr}
\usepackage{parskip}
\usepackage{microtype}
\usepackage[english, ukrainian]{babel}
\usepackage{url}
\usepackage{amsmath}
\usepackage{amssymb}
\usepackage{amsthm}
\usepackage{longtable}
\usepackage{booktabs}
\usepackage{hyperref}
\usepackage{titlesec}
\usepackage{tikz}
\usetikzlibrary{shapes.geometric, arrows, arrows.meta, positioning, fit, backgrounds}

\definecolor{nato}{HTML}{293757}
\definecolor{footergray}{gray}{0.65}

\geometry{a4paper,left=3cm,right=3cm,top=3cm,bottom=3cm}
\renewcommand{\familydefault}{\sfdefault}
\setcounter{secnumdepth}{3}

\DeclareFontFamily{T2A}{phv}{}
\DeclareFontShape{T2A}{phv}{m}{n}{<->ssub * cmss/m/n}{}
\DeclareFontShape{T2A}{phv}{bx}{n}{<->ssub * cmss/bx/n}{}
\DeclareFontShape{T2A}{phv}{m}{it}{<->ssub * cmss/m/it}{}
\DeclareFontShape{T2A}{phv}{bx}{it}{<->ssub * cmss/bx/it}{}

\titleformat{\section}{\color{nato}\Huge\bfseries}{\thesection.}{1em}{}
\titlespacing*{\section}{0pt}{30pt}{12pt}
\titleformat{\subsection}{\color{nato}\LARGE\bfseries}{\thesubsection.}{1em}{}
\titlespacing*{\subsection}{0pt}{20pt}{8pt}
\titleformat{\subsubsection}{\color{nato}\Large\bfseries}{\thesubsubsection.}{1em}{}
\titlespacing*{\subsubsection}{0pt}{10pt}{6pt}
\titleformat{\paragraph}{\color{nato}\large\bfseries}{\theparagraph.}{1em}{}
\titlespacing*{\paragraph}{0pt}{8pt}{4pt}

\fancyhf{}
\fancyhead[R]{\small\color{footergray} 19 серпня 2026}
\fancyfoot[L]{\small\color{footergray} ERP/1. Операційна система}
\fancyfoot[R]{\small\color{footergray}\thepage}
\renewcommand{\headrulewidth}{0pt}
\renewcommand{\footrulewidth}{0.4pt}
\renewcommand{\footrule}{\color{footergray}\hrule width \textwidth height 0.4pt}

\newcommand\HUGE[1]{\fontsize{40}{40}\selectfont #1}

\let\originalfootnotesize\footnotesize
\let\oldverbatim\verbatim
\let\oldendverbatim\endverbatim
\renewenvironment{verbatim}{\originalfootnotesize\oldverbatim}{\oldendverbatim}

\title{\textbf{OS.1}: Frozen Baseline Declaration \\ Erlang/OTP~20 + LibreSSL~3.1.5}
\author{Namdak Tonpa (maxim@synrc.com), Synrc Research}
\date{\today}

\hypersetup{
    colorlinks=true,
    linkcolor=blue,
    filecolor=magenta,
    urlcolor=cyan,
    pdftitle={OS.1 Frozen Baseline: OTP 20 + LibreSSL 3.1.5},
    pdfpagemode=FullScreen,
}

\begin{document}

%----------------------------------------------------------------------
\begin{titlepage}
  \AddToShipoutPictureBG*{\AtPageLowerLeft{\color{nato}\rule{\paperwidth}{\paperheight}}}
  \color{white}\thispagestyle{empty}\vspace*{3cm}\centering
  {\Huge \textbf{Zen Crypted Operating System} \par} \vspace{1.5cm}
  {\HUGE \textbf{Технічна декларація} \par} \vspace{1.5cm}
  {\LARGE \textbf{Заморожений базовий рівень \\ Erlang/OTP~20 з LibreSSL~3.1.5 \\ для контрольованого подальшого патчингу} \par} \vspace{1.5cm}
  {\large На основі вимог високого рівня довіри \par}
  {\large Для статичних Protection Domain seL4 \par}
  {\large Формалізовано згідно з практиками ISO/IEC/IEEE 29148 \par} \vfill
  {\large 2026 \copyright\ Криптографічні Телесистеми \par}
\end{titlepage}

\maketitle
\tableofcontents

%----------------------------------------------------------------------
\begin{abstract}
This document formally declares a frozen software baseline consisting of Erlang/OTP~20
(final maintenance line) statically linked against LibreSSL~3.1.5. The freeze is instituted
to enable rigorous, reviewable and reproducible further patching in high-assurance environments,
particularly those employing the seL4 microkernel with static Protection Domains, single-core
execution and a minimised cryptographic surface (including the synrc/ca pure-ASN.1 stack).
The declaration establishes the precise component versions, the rationale for the freeze,
the permitted scope of subsequent patches, the exact \texttt{:crypto} and \texttt{:public\_key}
call surface, and the governance principles under which any modifications may be introduced.
\end{abstract}

%----------------------------------------------------------------------
\section{Introduction}
\label{sec:intro}

Modern cryptographic and runtime ecosystems evolve rapidly. Continuous upstream updates,
while beneficial for general-purpose systems, introduce unacceptable variability and review
burden for high-assurance, formally constrained or long-lived embedded deployments.
In such contexts a deliberately frozen and fully characterised baseline is a prerequisite
for meaningful security evaluation, deterministic behaviour and controlled remediation.

This paper declares the following frozen baseline:

\begin{itemize}
  \item \textbf{Erlang/OTP~20} (final maintenance releases, including 20.3.8.$x$);
  \item \textbf{LibreSSL~3.1.5} (December 2020), used as the sole cryptographic provider
        via the OTP \texttt{crypto} NIF, preferably under static linkage.
\end{itemize}

The baseline is intended for systems that prioritise auditability, minimality and isolation
(e.g.\ seL4 static Protection Domains) over access to the newest cryptographic primitives
or protocol features.

%----------------------------------------------------------------------
\section{Rationale for the Freeze}
\label{sec:rationale}

\begin{enumerate}
  \item \textbf{Reviewability.}  
        OTP~20 and LibreSSL~3.1.5 constitute a finite, historically bounded codebase.
        Their combined surface can be subjected to exhaustive manual and automated review.

  \item \textbf{Reproducibility.}  
        Exact source trees, build configurations and linkage methods can be archived and
        regenerated years later.

  \item \textbf{Compatibility with constrained execution environments.}  
        Single-core, non-SMP, statically linked builds of OTP~20 map cleanly onto seL4
        Protection Domains with minimal capability sets.

  \item \textbf{Controlled patching surface.}  
        By freezing the major versions, all subsequent security or correctness fixes become
        explicit, localised patches against a known baseline.

  \item \textbf{Alignment with existing high-assurance stacks.}  
        The baseline is compatible with pure-Erlang ASN.1 approaches (exemplified by synrc/ca)
        that minimise reliance on high-level LibreSSL X.509 APIs while still requiring a
        stable cryptographic substrate.
\end{enumerate}

%----------------------------------------------------------------------
\section{Precise Baseline Definition}
\label{sec:baseline}

\begin{table}[ht]
\centering
\begin{tabular}{@{}llp{6.2cm}@{}}
\toprule
\textbf{Component} & \textbf{Version / Identifier} & \textbf{Notes} \\
\midrule
Erlang/OTP              & 20.3.8.$x$ (final)          & SMP disabled or strictly limited; static NIFs preferred \\
\texttt{crypto}         & OTP~20 corresponding       & Sole interface to LibreSSL \\
\texttt{public\_key}    & OTP~20 corresponding       & Retained for records and pkix helpers \\
\texttt{asn1}           & OTP~20 corresponding       & Required by public\_key and synrc ASN.1 modules \\
LibreSSL                & 3.1.5                      & Static linking; unnecessary features disabled \\
Linkage                 & Static (preferred)         & \texttt{libcrypto.a} linked into the crypto NIF \\
Higher-level stack      & synrc/ca (or equivalent)   & Pure-ASN.1 certificate authority layer \\
\bottomrule
\end{tabular}
\caption{Frozen baseline components.}
\label{tab:baseline}
\end{table}

Build flags for LibreSSL~3.1.5 shall exclude shared libraries, engines, deprecated protocols
and non-essential ciphersuites where feasible. OTP shall be configured to avoid dynamic
loading of the SSL library.

%----------------------------------------------------------------------
\section{Scope of the Freeze}
\label{sec:scope}

\paragraph{Frozen (immutable without formal amendment)}
\begin{itemize}
  \item Major and minor version of Erlang/OTP (20.$x$);
  \item Exact LibreSSL release (3.1.5);
  \item Core architectural decisions (single-core preference, static linkage, retention of \texttt{public\_key}).
\end{itemize}

\paragraph{Permitted under controlled patching}
\begin{itemize}
  \item Security fixes addressing known CVEs or demonstrated vulnerabilities;
  \item Correctness patches required for deterministic behaviour on the target platform;
  \item Minimal adaptations required for seL4 integration or static Protection Domain constraints;
  \item Removal or disabling of unused code paths (attack-surface reduction);
  \item Documentation and build-system improvements that do not alter runtime semantics.
\end{itemize}

Any patch must be explicitly documented with rationale and risk analysis, applied against
the frozen baseline, and subject to independent review before integration.

%----------------------------------------------------------------------
\section{Cryptographic Call Surface}
\label{sec:crypto-surface}

The following tables enumerate the exact entry points into the cryptographic substrate
that are exercised by the synrc/ca + \texttt{x509} stack when running on the frozen baseline.

\subsection{Direct / Primary \texttt{:crypto} Calls}
\label{sec:crypto-direct}

\begin{longtable}{@{}p{4.1cm}p{4.3cm}p{3.2cm}p{3.2cm}@{}}
\toprule
\textbf{Function} & \textbf{Typical form} & \textbf{Used for} & \textbf{Called from} \\
\midrule
\endfirsthead
\toprule
\textbf{Function} & \textbf{Typical form} & \textbf{Used for} & \textbf{Called from} \\
\midrule
\endhead
\bottomrule
\endfoot

\texttt{:crypto.generate\_key/2} &
\texttt{generate\_key(:ecdh, Curve)} or via \texttt{public\_key} &
EC key pair generation &
\texttt{X509.PrivateKey.new\_ec/1} $\to$ \texttt{CA.CSR.root/2}, \texttt{server/2}, \texttt{client/2} \\

\texttt{:crypto.generate\_key/3} &
\texttt{generate\_key(:eddsa, Curve, Priv)} &
EdDSA public derivation &
\texttt{x509} (Ed25519/Ed448; secondary for secp* focus) \\

\texttt{:crypto.sign/4} &
\texttt{sign(Algorithm, DigestType, Msg, Key)} &
ECDSA / RSA signatures &
CSR \& certificate signing (via \texttt{x509}/\texttt{public\_key}) \\

\texttt{:crypto.verify/5} &
\texttt{verify(Algorithm, DigestType, Msg, Sig, Key)} &
Signature verification &
\texttt{X509.CSR.valid?/1}, certificate checks \\

\texttt{:crypto.strong\_rand\_bytes/1} &
\texttt{strong\_rand\_bytes(N)} &
Serial numbers, IVs, nonces, salt &
Certificate serials, AES IVs, key encryption \\

\texttt{:crypto.compute\_key/4} &
\texttt{compute\_key(:ecdh, OthersPub, MyPriv, Curve)} &
ECDH shared secret &
\texttt{CA.AES} shared-secret helpers \\

\texttt{:crypto.crypto\_one\_time/4--5} &
\texttt{crypto\_one\_time(Cipher, Key, IV, Data, Flag)} &
AES-256-ECB/CBC/GCM/CCM &
\texttt{CA.AES} encrypt/decrypt \\

\texttt{:crypto.hash/2} &
\texttt{hash(sha256 $\mid$ sha384 $\mid$ \ldots, Data)} &
Digests for signatures / HKDF &
Indirect via \texttt{public\_key}/\texttt{x509} and \texttt{CA.HKDF}/\texttt{CA.KDF} \\

\texttt{:crypto.hmac/3} &
HMAC (OTP~20 equivalent) &
HKDF / key derivation &
\texttt{CA.HKDF} / \texttt{CA.KDF} paths \\

\texttt{:crypto.privkey\_to\_pubkey/2} &
(engine / modern path) &
Derive public from private &
\texttt{X509.PublicKey.derive/1} (engine case) \\
\end{longtable}

\subsection{Calls Mediated by \texttt{:public\_key}}
\label{sec:public-key}

These constitute the most common path inside synrc/ca.

\begin{longtable}{@{}p{5.2cm}p{3.8cm}p{5.5cm}@{}}
\toprule
\textbf{\texttt{:public\_key} function} & \textbf{Ultimately uses} & \textbf{Purpose in synrc/ca + x509} \\
\midrule
\endfirsthead
\toprule
\textbf{\texttt{:public\_key} function} & \textbf{Ultimately uses} & \textbf{Purpose in synrc/ca + x509} \\
\midrule
\endhead
\bottomrule
\endfoot

\texttt{public\_key:generate\_key/1} &
\texttt{:crypto.generate\_key} &
\texttt{X509.PrivateKey.new\_ec/1}, \texttt{new\_rsa/1} \\

\texttt{public\_key:sign/3} or \texttt{/4} &
\texttt{:crypto.sign} &
Certificate \& CSR signing \\

\texttt{public\_key:verify/4} or \texttt{/5} &
\texttt{:crypto.verify} &
CSR / signature validation \\

\texttt{public\_key:pkix\_sign/2} &
\texttt{:crypto.sign} &
Older-style certificate signing \\

\texttt{public\_key:pem\_encode/1}, \texttt{pem\_decode/1} &
mostly pure + crypto for encrypted PEM &
Key \& certificate PEM I/O in \texttt{CA.CSR} \\

\texttt{public\_key:pkix\_encode/3}, \texttt{pkix\_decode\_cert/2} &
ASN.1 + crypto for signatures &
DER handling in CMP/EST replies \\

\texttt{public\_key:der\_encode/2}, \texttt{der\_decode/2} &
ASN.1 only &
Structure manipulation (no direct crypto) \\
\end{longtable}

\subsection{Recommended Minimal Whitelist}
\label{sec:whitelist}

For the smallest practical attack surface that still supports the full EC-focused
synrc/ca functionality on the frozen baseline, the following entry points are sufficient:

\begin{verbatim}
:crypto.generate_key/2
:crypto.sign/4
:crypto.verify/5
:crypto.strong_rand_bytes/1
:crypto.hash/2
:crypto.compute_key/4          % only if CA.AES / ECDH retained
:crypto.crypto_one_time/5      % only if AES support retained
\end{verbatim}

All higher-level certificate structure handling remains in pure Erlang ASN.1 modules
and does not enlarge the cryptographic TCB.

%----------------------------------------------------------------------
\section{Patching Policy and Governance}
\label{sec:policy}

\begin{enumerate}
  \item \textbf{No silent upgrades.} Version numbers of the baseline components remain fixed.
        Patches are recorded as delta sets against the declared baseline.

  \item \textbf{Whitelisting of cryptographic entry points.} The \texttt{:crypto} NIF surface
        should be restricted to the minimal set listed in Section~\ref{sec:whitelist}.

  \item \textbf{Traceability.} Every binary artefact must be accompanied by exact source hashes,
        full build configuration, list of applied patches, and a mapping of high-level operations
        to LibreSSL functions.

  \item \textbf{Exception process.} Introduction of any newer cryptographic algorithm or protocol
        feature requires a formal amendment to this declaration and a re-evaluation of the isolation
        and review arguments.
\end{enumerate}

%----------------------------------------------------------------------
\section{Security Considerations and Residual Risk}
\label{sec:risk}

LibreSSL~3.1.5 post-dates the Heartbleed vulnerability and removes a substantial amount of
legacy OpenSSL complexity. Nevertheless it is a 2020-era library and therefore carries residual
risk relative to contemporary cryptographic practice. The freeze accepts this risk in exchange
for a stable, reviewable foundation.

Mitigations inherent in the intended deployment model include:

\begin{itemize}
  \item Strong isolation via seL4 static Protection Domains;
  \item Aggressive reduction of network and protocol attack surface;
  \item Preference for pure-ASN.1 certificate handling (synrc/ca style);
  \item Capability minimisation and single-core execution.
\end{itemize}

The baseline does \emph{not} claim to provide state-of-the-art cryptographic agility.
It claims only to provide a well-defined, frozen and patchable foundation upon which
higher-assurance arguments can be constructed.

%----------------------------------------------------------------------
\section{Conclusion}
\label{sec:conclusion}

We hereby declare Erlang/OTP~20 together with LibreSSL~3.1.5 as a frozen baseline for further
controlled patching. This declaration prioritises reviewability, reproducibility and compatibility
with high-assurance isolation mechanisms over continuous upstream evolution. All subsequent work
on systems adopting this baseline shall treat the stated versions as immutable reference points,
with any modifications introduced solely through explicit, reviewed patches.

This freeze is intentional, documented and subject to formal governance. It constitutes a deliberate
engineering decision in favour of long-term inspectability and deterministic behaviour.

%----------------------------------------------------------------------
\section*{References (indicative)}

\begin{itemize}
  \item Erlang/OTP~20 release and maintenance documentation.
  \item LibreSSL~3.1.5 release notes and change log.
  \item seL4 Reference Manual and Protection Domain model.
  \item synrc/ca technical description and ASN.1 module set.
  \item Analyses of OTP \texttt{crypto} NIF surfaces and LibreSSL linkage practices.
\end{itemize}

\end{document}
